\documentclass[a4paper,11pt]{article}
\usepackage[utf8]{inputenc}
\usepackage[T1]{fontenc}
\usepackage[french]{babel}
\usepackage[left=.7cm,right=.7cm,top=.5cm,bottom=.2cm]{geometry}
\usepackage{enumitem}
\usepackage{hyperref}
\usepackage{xcolor}
\usepackage{titlesec}
\usepackage{fontawesome5}
\usepackage{lmodern}
\usepackage[normalem]{ulem}

% --- Couleurs Crédit Agricole (Vert institutionnel) ---
\definecolor{primary}{RGB}{0, 102, 78} % Vert Crédit Agricole
\definecolor{secondary}{RGB}{60, 60, 60}

% --- Configuration des liens ---
\hypersetup{
    colorlinks=true,
    linkcolor=primary,
    filecolor=primary,
    urlcolor=primary,
}

% --- Formatage des titres ---
\titleformat{\section}{\normalsize\bfseries\color{primary}\uppercase}{}{0em}{}[\titlerule]
\titlespacing{\section}{0pt}{8pt}{5pt}

% --- Commandes personnalisées ---
\newcommand{\resumeItem}[1]{\item\small{#1}}
\newcommand{\resumeSubheading}[4]{
  \vspace{-2pt}\item
    \begin{tabular*}{0.97\textwidth}[t]{l@{\extracolsep{\fill}}r}
      \textbf{#1} & \textit{\small #2} \\
      \textit{\small #3} & \textit{\small #4} \\
    \end{tabular*}\vspace{-5pt}
}

\begin{document}

% --- EN-TÊTE ---
\begin{center}
    {\large \textbf{Loïc Aron MBASSI EWOLO}} \\ \vspace{3pt}
    \textbf{\large Stage : Assistant Data \& AI Engineer – Copilot Agentique d'Assistance au Développement} \\
    \textit{\small Disponible dès avril 2026 pour 4 à 6 mois} \\ \vspace{3pt}
    \small
    \faMapMarker\ Cergy, France (Mobile Île-de-France) \quad
    \faPhone\ (+33) 7 59 52 33 98 \quad
    \faEnvelope\ \href{mailto:loicaron.mbassiewolo@ensea.fr}{loicaron.mbassiewolo@ensea.fr} \\ \vspace{2pt}
    \faLinkedin\ \href{https://www.linkedin.com/in/loic-aron-mbassi-ewolo-3942a9246/}{linkedin.com} \quad
    \faGithub\ \href{https://github.com/Nameless0l}{Nameless0l} \quad
    \faGlobe\ \href{https://mbassiloic.tech}{mbassiloic.tech}
\end{center}

\vspace{-5pt}

% --- PROFIL ---
\noindent\small{Élève-ingénieur en double diplôme à l'\textbf{ENSEA} avec une formation riche et variée en \textbf{Mathématiques, Électronique, Informatique et Programmation}. 
J'aime explorer de nouvelles approches et technologies ainsi que relever des défis open source. 
Ce dont témoignent la note de \textbf{19,25/20} en physique obtenue au Baccalauréat S, la 2\textsuperscript{e} place au hackathon IndabaX \textbf{Africa} et mes divers travaux en \textbf{systèmes multi-agents} (Google ADK, JADE) et \textbf{génération de code assistée par LLM}.}

\vspace{.5pt}

% --- FORMATION ---
\section{Formation}
\begin{itemize}[leftmargin=*, label={.}, nosep]
    \item \textbf{Double diplôme : École Nationale Supérieure de l'Électronique et de ses Applications (ENSEA)} \hfill \textit{2025 -- Présent} \\
    \small{2ème année - Spécialité Intelligence Artificielle et Traitement du Signal, Ingénieur de conception}
    \vspace{2pt}
    \item \textbf{École Nationale Supérieure Polytechnique de Yaoundé — Génie Informatique} \hfill \textit{2021 -- Présent} \\
    \small{5ème année - Spécialité Génie Informatique, Ingénieur de conception}
\end{itemize}

% --- COMPÉTENCES TECHNIQUES ---
\section{Compétences Techniques}
\begin{itemize}[leftmargin=*, nosep]
    \item \small{\textbf{Agents IA \& LLM :} Google ADK, JADE (Java), LangChain, RAG, Prompt Engineering, API Gemini/OpenAI, Fine-tuning}
    \item \small{\textbf{Développement :} Python (avancé), Java, TypeScript, PHP | \textbf{Frameworks :} FastAPI, Spring Boot, Laravel, Next.js}
    \item \small{\textbf{DevOps \& CI/CD :} Git/GitHub, Docker, GitHub Actions, tests unitaires \& intégration, revue de code}
    \item \small{\textbf{IA/ML :} PyTorch, TensorFlow, Hugging Face Transformers, NLP, Computer Vision}
    \item \small{\textbf{Méthodologies :} Agile/Scrum, TDD, Architecture Microservices, Documentation technique}
\end{itemize}

% --- PROJETS AGENTS IA & GÉNÉRATION DE CODE ---
\section{Projets – Agents IA \& Automatisation du Développement}
\begin{itemize}[leftmargin=*, label={}]

    \item \textbf{Laravel API Generator, Open Source} \hfill \textit{2024 -- Présent, personnel}
    \vspace{-2pt}
    \begin{itemize}[leftmargin=0.4cm, nosep, label=\tiny$\bullet$]
        \item \small{Conception d'un package assisté par IA pour la \textbf{génération de projets} (à partir de diagrammes UML/AUML) en suivant les bonnes pratiques}
        \item \small{Parsing et analyse de fichiers XML, architecture modulaire, intégration de \textbf{LLM} pour génération et assistance au code}
        \item \small{\textbf{Compétences :} Python pour métaprogrammation, tests unitaires automatisés, documentation automatique — \href{https://github.com/Nameless0l}{\faGithub}}
    \end{itemize}
    \vspace{4pt}

    \item \textbf{Système Multi-Agents avec Google ADK} \hfill \textit{2024 -- 2025, personnel}
    \vspace{-2pt}
    \begin{itemize}[leftmargin=0.4cm, nosep, label=\tiny$\bullet$]
        \item \small{Développement d'un système d'agents IA collaboratifs orchestrant \textbf{5 agents spécialisés} pour recommandations agricoles personnalisées}
        \item \small{Orchestration multi-agents avec Google ADK, \textbf{RAG} (Retrieval-Augmented Generation), intégration API Gemini}
        \item \small{\textbf{Architecture :} Coordination automatique inter-agents avec résolution de conflits, vectorisation LangChain, Poetry, FastAPI, Docker}
    \end{itemize}
\end{itemize}

% --- EXPÉRIENCES PROFESSIONNELLES ---
\section{Expériences Professionnelles}
\begin{itemize}[leftmargin=*, label={}]

    \resumeSubheading
      {Stage de Recherche à distance chez MILA}{Juin -- Août 2025}
      {Institut Québécois d'Intelligence Artificielle}{Québec, Canada}
      \resumeItem{Grokking (de OpenAI), Généralisation des réseaux de neurones après sur-apprentissage}
      \begin{itemize}[leftmargin=0.4cm, nosep, label=\tiny$\bullet$]
        \item \small{Reproduction d'expériences SOTA pour analyser le comportement des modèles face au surapprentissage}
        \item \small{Implémentation de pipelines de tests statistiques sous \textbf{PyTorch} pour valider les hypothèses expérimentales}
      \end{itemize}
    \vspace{2pt}

    \resumeSubheading
      {Développeur Full-Stack (Fintech) — CSC}{Oct. 2024 -- Août 2025}
      {Douala, Cameroun}{}
      \resumeItem{Conception et développement d'une API fintech sécurisée et d'interfaces de gestion}
      \begin{itemize}[leftmargin=0.4cm, nosep, label=\tiny$\bullet$]
        \item \small{Architecture microservices, bases de données relationnelles, containerisation (Docker), \textbf{CI/CD}}
        \item \small{Tests unitaires \& d'intégration (TDD), revue de code, optimisation des performances — rigueur et qualité du code}
      \end{itemize}
    \vspace{2pt}

\end{itemize}

% --- PROJETS COMPLÉMENTAIRES ---
\section{Autres Projets Significatifs}
\begin{itemize}[leftmargin=*, label={}]

    \item \textbf{SOSDOCTA, Belgique} \hfill \textit{Fév. 2022 -- Nov. 2024, freelance}
    \vspace{-2pt}
    \begin{itemize}[leftmargin=0.4cm, nosep, label=\tiny$\bullet$]
        \item \small{Application de consultation et gestion des patients avec télémédecine intégrée}
        \item \small{Laravel, API REST sécurisées (OAuth2/JWT), Docker Compose, maintenance 2+ ans}
    \end{itemize}
    \vspace{2pt}

    \item \textbf{Upgrade Jetson - Challenge CoHoMa} \hfill \textit{Présent, académique}
    \vspace{-2pt}
    \begin{itemize}[leftmargin=0.4cm, nosep, label=\tiny$\bullet$]
        \item \small{Mise en fonction des Jetsons pour des robots de l'armée, intégration et optimisation des modèles d'IA légers (YOLOv10)}
        \item \small{Déploiement Docker sur systèmes embarqués, optimisation TensorFlow/PyTorch, Linux embarqué, navigation SLAM}
    \end{itemize}
    \vspace{2pt}

    \item \textbf{GoFind - Plateforme Microservices Multi-Services} \hfill \textit{2024, académique}
    \vspace{-2pt}
    \begin{itemize}[leftmargin=0.4cm, nosep, label=\tiny$\bullet$]
        \item \small{Architecture microservices complète : Laravel, NestJS, MongoDB, communication inter-services avec \textbf{RabbitMQ}}
    \end{itemize}
    \vspace{2pt}

    \item \textbf{Système de simulation taxi multiprocessus en C — PROJET-TAXI} \hfill \textit{2023 -- 2024}
    \vspace{-2pt}
    \begin{itemize}[leftmargin=0.4cm, nosep, label=\tiny$\bullet$]
        \item \small{Architecture distribuée simulant clients, motos-taxis et centre opérationnel : IPC POSIX, signaux UNIX, mémoire partagée}
        \item \small{\textbf{Compétences :} programmation système bas-niveau (C), synchronisation, Makefiles modulaires — \href{https://github.com/Nameless0l/PROJET-TAXI}{\faGithub}}
    \end{itemize}

\end{itemize}

% --- ACTIVITÉS & LEADERSHIP ---
\section{Activités Académiques \& Leadership}
\begin{itemize}[leftmargin=*, nosep]
    \item \small{\textbf{Chef cellule Projet, AI Cell ENSPY} \hfill \textit{2024 -- Présent} | Coordination de projets IA, workshops ML, mentorat}
    \item \small{\textbf{Hackathons :} 2\textsuperscript{e} place IndabaX Africa (2025) | 1\textsuperscript{re} place CITS National \& Mountain Hub Regional (2024)}
    \item \small{\textbf{Workshops animés :} Laravel, React.js, Machine Learning, Fine-tuning de modèles de langage}
\end{itemize}

% --- CERTIFICATIONS ---
\section{Certifications}
\begin{itemize}[leftmargin=*, nosep]
    \item \small{Supervised Machine Learning (DeepLearning.AI) | Python for Data Science \& Development (IBM) | Agents IA (IBM, en cours)}
    \item \small{\textbf{Langues :} Français (Langue maternelle), Anglais (Communication scientifique)}
\end{itemize}

\end{document}
